\chapter{Introduction}

\section{The Access to Justice Problem in India}

Access to justice in India is not only a legal ideal but also a large-scale administrative challenge. Court backlogs remain extremely high across the judicial system, which means that delay itself becomes a serious barrier to justice. When cases take years to move, litigation becomes more expensive, evidence becomes harder to trace, and ordinary citizens are less able to sustain meaningful participation in the legal process.

This does not mean that courts should be automated. Judicial decision-making involves interpretation, justification, and public accountability in ways that no prediction system can replace. A more realistic goal is decision support: computational tools that help organise large volumes of case material, surface patterns across past disputes, and assist legal research or case preparation. Legal Judgment Prediction (LJP) sits within that narrower and more defensible ambition.

The Indian setting makes this especially important. Judgments are often long, noisy, citation-heavy, and spread across multiple hearings. As a result, the challenge is not just to classify a document, but to build a faithful representation of a dispute from complex judicial text. This thesis addresses that challenge through structured case modelling rather than flat document classification.

\section{What is Legal Judgement Prediction?}

Legal Judgment Prediction is the task of estimating some aspect of a court outcome from information contained in the case record \cite{aletras2016echr, katz2017supreme, medvedeva2020echr}. Different papers define the target differently: some predict violation versus non-violation, some predict appeal outcomes, and others predict charges or penalties. In all cases, the core idea is the same: given a representation of a dispute, can a model forecast the judicial result?

In this thesis, LJP is defined narrowly as outcome prediction over Indian court judgments after removing explicit decision-bearing text. The aim is not to generate legal reasoning or normative conclusions, but to test whether pre-decision signals --- facts, arguments, statutes, provisions, and precedents --- can support outcome prediction when represented in a structured way.

This definition matters because much of the LJP literature has been criticised for hidden leakage \cite{medvedeva2023doitright}. If a model can read operative text, ruling sections, or highly outcome-correlated metadata, then the task becomes answer recovery rather than prediction. For that reason, LJP is as much an experimental-design problem as it is a modelling problem. The central question is not only \emph{what} to predict, but also \emph{what information the model is allowed to use}. 

\section{Why LJP is Hard}

LJP is difficult because legal judgments are not ordinary classification inputs. They are long, internally heterogeneous documents in which facts, procedural history, lawyer submissions, precedent citations, and judicial reasoning are interwoven across many pages. Standard text models often truncate or simplify this structure, which risks losing legally decisive context \cite{chalkidis2020legalbert, mamakas2022longlegal}.

The problem is harder in India because the text is noisy and structurally inconsistent. Judgments contain abbreviations, citation variations, OCR errors, transliterated names, and multilingual or mixed-register language. Core legal entities such as statutes, provisions, courts, and parties rarely appear in one clean canonical form. On top of that, one dispute may generate several hearing-stage documents rather than a single final text. If these are treated as separate cases, both the input representation and the labels become unreliable.

A further difficulty is that legal reasoning is inherently relational. Outcomes depend not just on isolated sentences, but on how facts, arguments, statutes, provisions, precedents, and parties connect. At the same time, LJP is highly vulnerable to leakage. Judgments often contain phrases such as \emph{appeal allowed} or \emph{petition dismissed}, and even identity-heavy metadata may act as shortcuts. Finally, the task is cross-domain: different legal areas have different statutes, document styles, and outcome distributions. Together, these challenges make LJP a representation and evaluation problem first, and only then a modelling problem.

\section{Why Graph-Based Approaches?}

A graph-based approach is a natural fit for legal disputes because cases are structured collections of typed objects and relations. A case is heard in a court, involves parties and lawyers, cites statutes and precedents, and contains argument blocks linked to different sides. When all of this is flattened into one sequence, the model must recover structure that is already present in the source material. A graph preserves that structure explicitly.

Graphs are useful here for two reasons. First, they preserve heterogeneity: a statute, a party, and a facts section do not play the same legal role and should not be treated as identical inputs. Second, they allow cross-case reasoning through shared legal objects. If two cases cite the same statute or precedent, they are connected in a legally meaningful way, not merely through surface lexical similarity. This thesis therefore argues for a heterogeneous, leakage-audited, cross-case graph representation for Indian LJP.

\section{Contributions of This Work}

This thesis makes five main contributions.

\begin{enumerate}
    \item It constructs a large structured corpus of Indian court cases spanning five legal domains, with more than 73{,}000 final-case records after multi-hearing consolidation.

    \item It develops an end-to-end pipeline that converts raw judgments into leakage-audited, graph-ready representations using document collection, OpenNyAI enrichment, outcome labelling, multi-hearing merging, and explicit leakage filtering.

    \item It proposes a reasoning-focused heterogeneous knowledge graph schema in which cases are represented together with text blocks, parties, courts, statutes, provisions, and precedents, while legally meaningful sharing is separated from identity-based shortcuts.

    \item It applies a Heterogeneous Graph Transformer to this graph and evaluates it in the pooled cross-domain setting, showing that structured graph representations outperform text-only baselines on the thesis dataset.

    \item It complements prediction with graph analysis, ablations, and leakage-oriented reasoning, treating the graph not just as model input but as an empirical object that helps explain both performance and limitations.
\end{enumerate}

\section{Thesis Organisation}

The remainder of the thesis is organised as follows. Chapter 2 reviews prior work on legal judgment prediction, transformer-based legal NLP, graph methods, heterogeneous GNNs, and Indian legal NLP resources. Chapter 3 describes the full data pipeline, from document collection to OpenNyAI enrichment, outcome labelling, and multi-hearing merging. Chapter 4 presents the heterogeneous graph schema and the leakage-prevention strategy. Chapter 5 studies the graph structurally through statistics, bridge analysis, and per-domain patterns. Chapter 6 introduces the prediction model and reports the main results and ablations. Chapter 7 discusses implications, limitations, and future directions. Chapter 8 concludes the thesis.
